%=============================================================================!
% 15.2  SAMPLES THAT REMEMBER                                                 !
%=============================================================================!
\section{Samples that remember}
\label{sec:cv-correlation}

The samples of a run are not independent. In one step of
\SI{10}{\femto\second} an atom moves a small fraction of its distance to a
neighbour, so the potential energy, the kinetic energy and the pressure
change little from one step to the next, and a sample says much about the
one after it. This section measures that memory, finds how it enlarges the
error of a mean, and estimates it from a single run.

\subsection*{A series with memory}

A grain of dust in still air is pushed about by the molecules that strike
it, and its velocity changes over a time set by the friction of the air.
\Cref{sec:th-langevin} found the exact law of one component of such a
velocity over a step: the Ornstein-Uhlenbeck process of \cref{eq:th-ou}.
Measured in units of its own standard deviation, so that its variance is
1, it reads
\begin{equation}
   x_{k+1} = c\,x_k + \sqrt{1 - c^2}\;z_k , \qquad 0 \le c < 1 ,
   \label{eq:cv-ou}
\end{equation}
with $z_k$ a fresh Gaussian number of mean 0 and variance 1 at each step
(the $g$ of \cref{eq:th-ou}, renamed because $g$ is needed below).
The old value keeps the fraction $c$ of itself, and the noise supplies the
variance that this loses, $1 - c^2$. With $c = 0$ each sample is new; with
$c$ close to 1 the series wanders slowly. \Cref{fig:cv-ou}(a) draws 400
samples with $c = 0.95$, read as samples \SI{10}{\femto\second} apart,
against 400 independent ones: the first stays on one side of zero for up
to \SI{0.98}{\pico\second} at a time, while the second changes sign on
average every other sample.

\subsection*{The autocorrelation function}

The memory is measured by the covariance of two quantities,
\begin{equation*}
   \operatorname{cov}(x, y) = \ens{(x - \ens{x})(y - \ens{y})}
   = \ens{xy} - \ens{x}\ens{y} ,
\end{equation*}
the bracket of \cref{sec:sm-probability} that vanished for independent
quantities, which by the same expansion makes $\operatorname{var}(x + y) =
\operatorname{var}(x) + \operatorname{var}(y) + 2\operatorname{cov}(x, y)$.
A run that has settled has the same statistics wherever in it they are
measured, so the covariance of the samples $A_i$ and $A_{i+k}$ depends on
their separation $k$ alone. Divided by the variance, it is the
autocorrelation function
\begin{equation*}
   \operatorname{corr}(k) = \frac{\operatorname{cov}(A_i,
   A_{i+k})}{\operatorname{var}(A)} ,
\end{equation*}
equal to 1 at $k = 0$ and falling towards 0 as the series forgets. For the
series of \cref{eq:cv-ou}, multiplying both sides by $x_0$ and averaging
gives
\begin{equation*}
   \ens{x_{k+1}x_0} = c\ens{x_kx_0} + \sqrt{1 - c^2}\,\ens{z_kx_0}
   = c\ens{x_kx_0} ,
\end{equation*}
since $z_k$ is drawn after $x_0$ and independently of it, with mean zero.
Applying this $k$ times from $\ens{x_0x_0} = 1$ gives $\ens{x_kx_0} = c^k$,
and with the means zero and the variance 1 this is $\operatorname{corr}(k)
= c^k$, which falls by the factor $c$ each step. Here $\dt$ is the time
between samples, one integration step in this chapter's runs, and $c^k =
\mathrm{e}^{-k\dt/\tau_{\mathrm c}}$ with the correlation time $\tau_{\mathrm
c} = -\dt/\ln c$, which is $1/\gamma$ for Langevin's $c =
\mathrm{e}^{-\gamma\dt}$. For $c = 0.95$ and $\dt = \SI{10}{\femto
\second}$, $\tau_{\mathrm c} = \SI{195.0}{\femto\second}$.

\subsection*{The variance of a correlated mean}

The mean $\bar A$ of $n$ samples has the variance of their sum divided by
$n^2$. The square of a sum of deviations $\delta A_i = A_i - \ens{A}$ is
the sum of all their products, $(\sum_i\delta A_i)^2 =
\sum_i\sum_j\delta A_i\,\delta A_j$, so
\begin{equation*}
   \operatorname{var}(\bar A) = \frac{1}{n^2}\sum_{i=1}^n\sum_{j=1}^n
   \operatorname{cov}(A_i, A_j)
   = \frac{\sigma_A^2}{n^2}\sum_{i=1}^n\sum_{j=1}^n
   \operatorname{corr}(\lvert i - j\rvert) ,
\end{equation*}
the second form using the autocorrelation function at the lag $\lvert i -
j\rvert$, since $\operatorname{cov}(A_i, A_j) = \operatorname{cov}(A_j,
A_i)$. The pairs with $i = j$ number $n$, each contributing 1; those $k$ apart,
with $k$ from 1 to $n - 1$, number $n - k$ with $j$ after $i$ and as many
with $j$ before. Collecting them,
\begin{equation}
   \operatorname{var}(\bar A)
   = \frac{\sigma_A^2}{n^2}\Bigl[n + 2\sum_{k=1}^{n-1}(n - k)
   \operatorname{corr}(k)\Bigr]
   = \frac{\sigma_A^2}{n}\Bigl[1 + 2\sum_{k=1}^{n-1}\Bigl(1 - \frac kn\Bigr)
   \operatorname{corr}(k)\Bigr] .
   \label{eq:cv-var}
\end{equation}
For independent samples the bracket is 1 and this is \cref{eq:cv-mean}. A
run is normally far longer than its memory, so $k/n$ is negligible
wherever $\operatorname{corr}(k)$ is not, and the bracket settles to the
statistical inefficiency\cite{Chodera2007}
\begin{equation*}
   g = 1 + 2\sum_{k=1}^{\infty}\operatorname{corr}(k) , \qquad
   \operatorname{var}(\bar A) = \frac{g\,\sigma_A^2}{n} .
\end{equation*}
Correlated samples spread the mean $\sqrt g$ times as widely as the same
number of independent ones: $n$ of them are worth $n_{\mathrm{eff}} = n/g$
independent samples. For the series of \cref{eq:cv-ou} the sum is
geometric: the series of \cref{sec:th-langevin} gives $\sum_{k=0}^{m-1}c^k
= (1 - c^m)/(1 - c)$, which tends to $1/(1 - c)$ as $c^m$ vanishes, and
without its first term, 1, $\sum_{k\ge1}c^k = c/(1 - c)$, so
\begin{equation*}
   g = 1 + \frac{2c}{1 - c} = \frac{1 + c}{1 - c} ,
\end{equation*}
which is 39.0 for $c = 0.95$: a run of 10\,000 such samples is worth 256
independent ones. Writing the run's length as $t_{\mathrm f} = n\dt$ and
\begin{equation*}
   \tau_{\mathrm{int}} = \tfrac12g\,\dt , \qquad
   \operatorname{var}(\bar A) = \frac{\sigma_A^2}{t_{\mathrm f}/
   (2\tau_{\mathrm{int}})} ,
\end{equation*}
puts the result in times: a run holds one independent sample for every
two integrated correlation times $\tau_{\mathrm{int}}$, however often it
is sampled, provided the samples lie much closer together than
$\tau_{\mathrm{int}}$. For the series above $\tau_{\mathrm{int}} =
\SI{195.0}{\femto\second}$, its $\tau_{\mathrm c}$ to the four figures
shown: $\tau_{\mathrm{int}} = (\frac12 + \sum_{k\ge1}c^k)\dt$ is the area
under $\mathrm{e}^{-t/\tau_{\mathrm c}}$ counted in trapezia of width
$\dt$, whose exact area is $\tau_{\mathrm c}$.
\Cref{fig:cv-ou}(c) measures the spread of the mean, 4000 series over: for
$n = 1000$ it is 0.1903, against 0.1956 from \cref{eq:cv-var} and $1/\sqrt
n = 0.0316$ for independent samples. For runs not many times longer than
the memory the bracket of \cref{eq:cv-var} has not yet reached $g$: 8.5 at $n = 10$, 31.4
at $n = 100$.

\begin{figure}[htb]
\centering
\includegraphics{ch15_convergence/ou}
\caption{Samples that remember: the series of \cref{eq:cv-ou} with $c =
0.95$, samples \SI{10}{\femto\second} apart. (a) 400 samples (teal)
against 400 independent ones (grey). (b) Its autocorrelation function
estimated from 10\,000 samples, against the exact $c^k$ (dashed); the
estimate first reaches zero at the dotted line. (c) The spread of the mean
of $n$ samples, 4000 series over, against $1/\sqrt n$ for independent
samples (dotted) and the square root of \cref{eq:cv-var} (dashed).}
\label{fig:cv-ou}
\end{figure}

\subsection*{Estimating $g$ from one run}

A run does not know its autocorrelation function and must estimate it.
With the deviations now taken from the run's own mean, $\delta A_i = A_i -
\bar A$ in place of $A_i - \ens{A}$, which the run does not know, and
$s'^2 = \frac1n\sum_i\delta A_i^2$, the estimate at the lag $k$ averages the
$n - k$ products available,
\begin{equation*}
   \widehat{\operatorname{corr}}(k)
   = \frac{1}{(n - k)\,s'^2}\sum_{i=1}^{n-k}\delta A_i\,\delta A_{i+k} ,
\end{equation*}
the hat marking an estimate. Each value carries noise, and beyond the
decay the estimates scatter about zero (\cref{fig:cv-ou}(b)): by 0.044
over the lags 300 to 1000 of 10\,000 samples of the series above, more
than $1/\sqrt n = 0.010$ because the products are themselves correlated.
Summed over every lag with the weights $1 - k/n$, the estimates give zero
in place of $g$, whatever the series. Multiplying out,
$(1 - k/n)\,\widehat{\operatorname{corr}}(k) = \sum_i\delta A_i\,\delta
A_{i+k}/(n s'^2)$, so
\begin{equation*}
   1 + 2\sum_{k=1}^{n-1}\Bigl(1 - \frac kn\Bigr)
   \widehat{\operatorname{corr}}(k)
   = \frac{\sum_i\delta A_i^2 + 2\sum_{k\ge1}\sum_i\delta A_i\,\delta
   A_{i+k}}{n s'^2}
   = \frac{\bigl(\sum_i\delta A_i\bigr)^2}{n s'^2} = 0 ,
\end{equation*}
the middle numerator being the square of the sum written out, and the
deviations from the run's own mean adding to zero. The sum must therefore
stop where the signal ends. The rule of Chodera and
co-workers\cite{Chodera2007} stops it where the estimate first reaches
zero; pymbar, and the listing below, look for that zero only beyond the
third lag. For the series of
\cref{fig:cv-ou}(b) that is the lag 119, \SI{1190}{\femto\second}, and $g$
comes out at 48.0 against the exact 39.0; over 200 such series the
estimate averages 41.6 and spreads by 8.7. From a run worth a few hundred
independent samples $g$ is uncertain by about a fifth, and the error bar,
which goes as $\sqrt g$, by about a tenth, since a small relative change
in $g$ changes $\sqrt g$ by half as much. The function
\texttt{statistical\_inefficiency} of \texttt{mdlab.analysis.stats}
applies the rule, with the autocorrelation function computed by the fast
Fourier transform for speed (a tool Chapter~16 explains); in it
\texttt{c} holds the estimates $\widehat{\operatorname{corr}}(k)$ at the
lags \texttt{k}:

\begin{code}
def statistical_inefficiency(series: ArrayLike, mintime: int = 3) -> float:
    a = np.asarray(series, dtype=float)
    n = len(a)
    if np.all(a == a[0]):
        raise ValueError("a constant series has no fluctuations")
    c = autocorrelation(a)[1 : n - 1]
    k = np.arange(1, n - 1)
    stop = np.nonzero((c <= 0) & (k > mintime))[0]
    end = stop[0] if len(stop) else len(c)
    g = 1 + 2 * np.sum(c[:end] * (1 - k[:end] / n))
    return float(max(g, 1.0))
\end{code}

\noindent It agrees with the same estimate in the \texttt{timeseries}
module of the pymbar package to nine figures (\texttt{test\_stats.py}).

\subsection*{The liquid}

The liquid of Chapter~13, followed for \SI{2}{\nano\second} under CSVR
with $\tau_{\mathrm T} = \SI{1}{\pico\second}$ and sampled every step,
shows two kinds of memory (\cref{fig:cv-acf}(a)). The potential energy,
the kinetic energy and the pressure lose part of their correlation within
\SI{0.1}{\pico\second}, as the atoms rattle against their neighbours:
$\operatorname{corr}$ is then 0.46, 0.63 and 0.66. The rest fades over
picoseconds, and sets $g$: 146 for $U$, 211 for $K$ and 188 for $P$, or
$\tau_{\mathrm{int}}$ of 0.73, 1.05 and \SI{0.94}{\pico\second}. The slow
part comes from the thermostat. At fixed energy, from the same liquid
(\cref{fig:cv-acf}(b)), $U$ and $K$, whose sum cannot change, have the
same function and $\tau_{\mathrm{int}} = \SI{77}{\femto\second}$, and $P$
has \SI{123}{\femto\second}. Under CSVR the total energy, which only the
thermostat changes, has $\tau_{\mathrm{int}} = \SI{1.71}{\pico\second}$,
close to the time $\tau_{\mathrm T}C_V/C_K = \SI{1.59}{\pico\second}$ over
which \cref{sec:th-berendsen} found a weak coupling to move the energy,
while the shear component $\mathsf{P}_{xy}$ of the pressure tensor, which
is not tied to the total energy, has \SI{0.18}{\pico\second} in both runs. The
length of run an average needs therefore depends on the thermostat as
well as on the liquid: a gentle coupling lets the energy wander slowly,
and every quantity tied to the energy inherits that slowness.

\begin{figure}[htb]
\centering
\includegraphics{ch15_convergence/acf}
\caption{The memory of liquid argon at \SI{135}{\kelvin}, from runs of
\SI{2}{\nano\second} sampled every \SI{10}{\femto\second}. (a) Under CSVR
with $\tau_{\mathrm T} = \SI{1}{\pico\second}$: the autocorrelation
functions of $U$, $K$, $P$ and the shear component $\mathsf{P}_{xy}$. (b)
The same at fixed energy, where $U$ and $K$ have the same function. (c)
The estimate of $g$ for $U$ summed up to each lag, under CSVR (teal) and
at fixed energy (ochre), with the lags at which the rule of this section
stops each sum (dotted).}
\label{fig:cv-acf}
\end{figure}

\Cref{fig:cv-acf}(c) shows the rule at work: under CSVR the running sum
for $U$ climbs until the estimate first reaches zero at \SI{6.16}{\pico
\second}, and at fixed energy it levels off before the sum stops at
\SI{0.70}{\pico\second}. Over the whole run under CSVR, the potential
energy of the 256 atoms is $U = (-12.8236 \pm
0.0047)$\,\si{\electronvolt}, $T = (134.81 \pm 0.22)$\,\si{\kelvin} and $P =
(0.0849 \pm 0.0004)$\,\si{\giga\pascal}. The error bars that leave out
$\sqrt g$, $s/\sqrt n$ over the 200\,001 samples, would be 12.1, 14.5 and
13.7 times smaller, and would claim a precision that the run does not
have.

\begin{keyidea}
Successive samples of a run are correlated, and $n$ of them spread the
mean as widely as $n/g$ independent ones, with the statistical
inefficiency $g = 1 + 2\sum_k\operatorname{corr}(k)$; a run of length
$t_{\mathrm f}$ holds $t_{\mathrm f}/(2\tau_{\mathrm{int}})$ independent
samples, $\tau_{\mathrm{int}} = g\dt/2$. The error of the mean is $\sqrt{g\,s^2/n}$, with $g$ estimated
from the run by summing its autocorrelation function until the estimate
first reaches zero.
\end{keyidea}

\notebook{15}{set the memory of a series and compare the naive and the
correct error bars}

\begin{selfcheck}
A series has $\operatorname{corr}(k) = 0.5^k$. What is $g$, and how many
independent samples are 1000 of its samples worth?
\end{selfcheck}
